\chapter{Introduction}
\label{ch:introduction}

\section{Motivation and problem}
\label{sec:motivation}

The topic came out of a discussion with my supervisor, Gregory Morse, before the
thesis registration. I was already working with AI agents during my internship
at Ericsson, and he suggested a topic in which agents would be used to discover
algorithms. After going through a few options I chose this one, because it
builds on what I do at work and takes it further: instead of using agents to
write ordinary code, the goal is to use them to find algorithms that are both
correct and efficient.

The problem is keeping a graph search structure up to date while the graph
changes. When an edge is inserted or deleted, a dynamic algorithm has to update
a BFS tree, a DFS tree or an LBFS ordering without recomputing it from scratch.
The idea is simple, but the details are easy to get wrong. Each structure has
invariants that must hold after every update: the BFS distances must agree with
the parent pointers, the DFS discover and finish times must nest properly, and
the LBFS ordering must keep its four-point property. One stale value or one
wrong tie-break is enough to break them. The known correct algorithms
\citep{demetrescu2005dynamic,even1981maintaining,friglioni2000,
yang2019dynamic,baswana2019dynamic,baswana2017,corneil2008unified} are
complicated enough that they are usually proved correct rather than tested.

I chose BFS, DFS and LBFS because they are the basic graph searches that many
other algorithms build on: BFS for shortest paths and layering, DFS for cycle
and connectivity tests, and LBFS for chordal-graph recognition and maximum
cardinality search \citep{morse2026bfs,brandstadt1997lexbfs,habib2000lexbfs,
berry2004mcs}. They also give a natural difficulty scale. Correct dynamic BFS
and DFS algorithms are already known \citep{yang2019dynamic,morse2026bfs}, so
they show whether the tool works at all; dynamic LBFS has been studied much
less, so it is the harder test. During development, the first problem I tried
was dynamic MCS, as a smaller starting point, before moving to these three.

Language models can now write and fix code \citep{chen2021codex,openai2025codex,
jimenez2024swebench}, and systems such as FunSearch and AlphaEvolve pair a model
with an evolutionary loop and have found new algorithms and mathematical
constructions \citep{funsearch2023,alphaevolve2025,alphatensor2022,lehman2023elm}.
The repair side has a longer history in genetic programming and automatic
program repair \citep{koza1992gp,legoues2012genprog}, and more recently in
agents that reflect on their own failures \citep{reflexion2023,
madaan2023selfrefine,yao2023react}. All of these systems judge a candidate by
running it on test inputs. For a dynamic graph algorithm that is not enough: a
candidate can pass every test and still fail on a small graph with one
particular order of updates \citep{claessen2000quickcheck,jia2011mutation}. I
saw this in my own evaluation. One candidate passed every graph up to four
vertices and failed 1 of 51{,}465 cases at five, and one of the seven one-line
mutations of the fully dynamic DFS of \citet{yang2019dynamic} fails only after
several updates, never after a single one (Section~\ref{sec:test-selftest}).

So I wanted a tool in which the agents and the checking are designed together.
Instead of testing on examples, the verifier checks a candidate on every graph
up to a size bound (one per isomorphism class), and the agents get the failing
case back as feedback. The tool can also be driven by an external coding agent,
not only by its own agents. To my knowledge, no existing tool does this for
dynamic graph algorithms. Within the stated bound, the tool reports that no
counterexample was found, and I use it to search for, verify and measure dynamic
BFS, DFS and LBFS algorithms, reporting both positive and negative results.

\section{Goal, scope, and success criteria}

The goal is a research tool that helps an algorithm designer develop a correct
and efficient dynamic graph algorithm: a Python package with a command line, a
local web interface and a harness mode for external coding agents, with an
experimental evaluation on a tiered set of problems.

The scope is a tiered problem set. Dynamic BFS and dynamic DFS form the
calibration tier, where a correct algorithm is known to be reachable, and
dynamic LBFS forms the open tier, where the search is not guaranteed to succeed.
The problems are directed except LBFS, which is undirected. The tiering is what
makes the evaluation interpretable: the calibration tier measures whether the
tool can find and verify a correct algorithm at all, and the open tier measures
whether it can do so on a problem where success is not assumed.

Two boundaries are stated explicitly as non-goals. First, verification is
bounded: a passing candidate has no counterexample within the stated graph-size
and sequence bound, and this is reported as bounded exhaustive verification,
never as a proof. Second, the sandbox that isolates untrusted candidate code
requires Linux with unprivileged user namespaces, so the tool is not portable to
other operating systems.

Success is measured by the three completion signals registered in the topic
form. The user interface can launch, monitor, and replay a multi-generation run
end to end. The verification engine validates known-good algorithms, including
the fully dynamic DFS of \citet{yang2019dynamic}, and rejects known-bad
mutations with zero false negatives on graphs up to five vertices; this
signal is measured on the committed fault corpus over the tested universe of
Section~\ref{sec:universe}, so it means that no seeded fault escaped, not that
every possible wrong program would be caught. And the
cooperative agents produce a correct dynamic BFS algorithm within a bounded
number of generations without human intervention.

\section{Three principal contributions}
\label{sec:contributions}

\begin{enumerate}
  \item \textbf{C1: an interactive research tool.} The difficult problem is
  observing and steering a live multi-agent search without importing the engine
  into the client. The solution is an event-sourced architecture: the engine
  appends to an immutable log, and the user interface is a pure reducer over
  that log, so the live view, the replay, and the status command all read the
  same source of truth. Evidence is the end-to-end browser run in
  Section~\ref{sec:test-ui} and the user documentation in
  Chapter~\ref{ch:user-documentation}.
  \item \textbf{C2: an isomorphism-aware verification engine.} The difficult
  problem is that dynamic outputs are not unique, so a candidate must be judged
  by an invariant or certificate rather than by equality, and the universe must
  cover every non-isomorphic graph up to the bound. The solution is a set of
  certificate checkers (BFS parent tree and distances; DFS laminarity with the
  absence of a forward-cross edge; LBFS four-point property) over a \texttt{nauty}
  enumeration, graded in two tiers. Evidence is the self-test, which rejects
  all 48 seeded faults within the tested universe, in Section~\ref{sec:test-selftest} and Table~\ref{tab:selftest}.
  \item \textbf{C3: cooperative agents with optional dynamic agent generation.}
  The difficult problem is turning a bounded counterexample into a targeted
  repair, and deciding when the existing agents have plateaued and a new
  specialised agent should be generated. The solution is an orchestrator with
  generator, reflector, and repairer roles, a structured reflection schema, and
  plateau detection that spawns and retires specialists. Evidence is the
  experimental evaluation in Section~\ref{sec:test-evaluation},
  Table~\ref{tab:eval1} and Table~\ref{tab:eval2}.
\end{enumerate}

Of the three, the verifier (C2) was the hardest to build. At the start I did not
have a clear picture of how it should work. Comparing a candidate's output with
a recomputed answer does not work, because a graph usually has many valid BFS or
DFS trees. The approach took shape over several consultations with my
supervisor, and I worked out the details with the help of AI agents. The agents
(C3) were the easiest part for me, because I already build AI agents at work and
had a clear idea of the roles from the beginning: one agent writes a candidate,
one explains why it failed, and one repairs it. The user interface (C1) was in
between: I had used React and Python web backends before, and I built most of
it with AI assistance (Section~\ref{sec:ai-use}).

\begin{table}[H]
  \centering
  \caption{Contribution-to-evidence map.}
  \label{tab:contribution-evidence}
  \begin{tabularx}{\textwidth}{@{}p{0.18\textwidth}X X@{}}
    \toprule
    \textbf{Contribution} & \textbf{What is original or difficult} & \textbf{Where it is shown} \\
    \midrule
    C1 & An event-sourced observation plane that steers a live multi-agent run and replays its generation history without importing the engine. & Section~\ref{sec:test-ui}, Figure~\ref{fig:ui-steer}, Figure~\ref{fig:ui-replay}, Chapter~\ref{ch:user-documentation}. \\
    \addlinespace[2pt]
    C2 & Judging non-unique dynamic outputs by certificate over an exhaustive \texttt{nauty} universe, with no seeded fault accepted within the tested universe. & Section~\ref{sec:test-selftest}, Table~\ref{tab:selftest}, Section~\ref{sec:universe}. \\
    \addlinespace[2pt]
    C3 & Structured failure reflection that turns a bounded counterexample into a targeted repair, and plateau-triggered dynamic agent generation. & Section~\ref{sec:test-evaluation}, Table~\ref{tab:eval1}, Table~\ref{tab:eval2}. \\
    \bottomrule
  \end{tabularx}
\end{table}

\section{Thesis roadmap}

This is a mixed research and engineering thesis: it builds a working tool and
uses it to produce experimental evidence, and it reports both the positive and
the negative results (the dynamic agent generation showed no significant effect,
Section~\ref{sec:test-evaluation}).

Chapter~\ref{ch:user-documentation} is the user manual,
Chapter~\ref{ch:requirements} the requirements baseline including the bounded
verification universe and the four refinements of the registered topic,
Chapter~\ref{ch:design} the design, Chapter~\ref{ch:implementation} the code
and its remaining shortcomings, Chapter~\ref{ch:testing} the tests and the
evaluation, and Chapter~\ref{ch:conclusion} the assessment against each
commitment. The appendices hold the reproducibility record
(Appendix~\ref{app:reproducibility}), the requirements catalogue
(Appendix~\ref{app:requirements}), the decision record
(Appendix~\ref{app:decisions}), the detailed tables
(Appendix~\ref{app:tables}) and the submission record
(Appendix~\ref{app:declarations}).

Generative AI tools were used in developing the software and in writing this
thesis. Section~\ref{sec:ai-use} gives the overview of their use, and the
generated figures are marked where they first appear.
